% =========================================================
% COPERTINA C - "Double beta"
% Spettro dell'energia sommata dei due elettroni per il 76Ge:
% continuo 2nubb (approssimazione di Primakoff-Rosen,
% Q = 2039 keV) e riga 0nubb a Q_bb (altezza esagerata,
% come nella Figure 0nu_vs_2nu del capitolo 3).
% =========================================================
\definecolor{coverBg}{HTML}{5a7692}
\definecolor{twoNu}{HTML}{F5A030}
\definecolor{zeroNu}{HTML}{E8434F}

\begin{titlepage}
\thispagestyle{empty}
\begin{tikzpicture}[remember picture, overlay]
  \fill[coverBg] (current page.south west) rectangle (current page.north east);
  \begin{scope}[shift={(current page.south west)}, x=1cm, y=1cm]
    \begin{scope}[yshift=6cm]
      % continuo 2nubb
      \fill[twoNu, fill opacity=0.22] (2.5,0) -- plot coordinates {(2.500,0.000) (2.605,0.194) (2.709,0.395) (2.814,0.604) (2.919,0.819) (3.023,1.039) (3.128,1.266) (3.233,1.496) (3.337,1.731) (3.442,1.970) (3.546,2.211) (3.651,2.454) (3.756,2.699) (3.860,2.945) (3.965,3.190) (4.070,3.436) (4.174,3.680) (4.279,3.923) (4.384,4.163) (4.488,4.400) (4.593,4.634) (4.698,4.863) (4.802,5.088) (4.907,5.308) (5.011,5.521) (5.116,5.729) (5.221,5.930) (5.325,6.123) (5.430,6.309) (5.535,6.487) (5.639,6.656) (5.744,6.816) (5.849,6.967) (5.953,7.109) (6.058,7.241) (6.163,7.363) (6.267,7.475) (6.372,7.576) (6.476,7.666) (6.581,7.746) (6.686,7.815) (6.790,7.873) (6.895,7.920) (7.000,7.957) (7.104,7.982) (7.209,7.996) (7.314,8.000) (7.418,7.993) (7.523,7.975) (7.628,7.947) (7.732,7.909) (7.837,7.861) (7.941,7.802) (8.046,7.735) (8.151,7.658) (8.255,7.572) (8.360,7.478) (8.465,7.375) (8.569,7.264) (8.674,7.146) (8.779,7.021) (8.883,6.889) (8.988,6.750) (9.093,6.606) (9.197,6.456) (9.302,6.301) (9.406,6.141) (9.511,5.978) (9.616,5.810) (9.720,5.640) (9.825,5.466) (9.930,5.291) (10.034,5.113) (10.139,4.934) (10.244,4.754) (10.348,4.574) (10.453,4.394) (10.558,4.214) (10.662,4.034) (10.767,3.856) (10.871,3.680) (10.976,3.505) (11.081,3.333) (11.185,3.163) (11.290,2.996) (11.395,2.833) (11.499,2.673) (11.604,2.517) (11.709,2.365) (11.813,2.217) (11.918,2.074) (12.023,1.935) (12.127,1.802) (12.232,1.673) (12.336,1.549) (12.441,1.431) (12.546,1.318) (12.650,1.211) (12.755,1.108) (12.860,1.011) (12.964,0.920) (13.069,0.834) (13.174,0.753) (13.278,0.677) (13.383,0.607) (13.488,0.541) (13.592,0.481) (13.697,0.425) (13.802,0.374) (13.906,0.327) (14.011,0.284) (14.115,0.245) (14.220,0.211) (14.325,0.180) (14.429,0.152) (14.534,0.128) (14.639,0.106) (14.743,0.087) (14.848,0.071) (14.953,0.057) (15.057,0.046) (15.162,0.036) (15.267,0.028) (15.371,0.021) (15.476,0.016) (15.580,0.011) (15.685,0.008) (15.790,0.006) (15.894,0.004) (15.999,0.002) (16.104,0.001) (16.208,0.001) (16.313,0.000) (16.418,0.000) (16.522,0.000) (16.627,0.000) (16.732,0.000) (16.836,0.000) (16.941,0.000) (17.045,0.000)} -- cycle;
      \draw[twoNu, line width=1.1pt, line join=round] plot coordinates {(2.500,0.000) (2.605,0.194) (2.709,0.395) (2.814,0.604) (2.919,0.819) (3.023,1.039) (3.128,1.266) (3.233,1.496) (3.337,1.731) (3.442,1.970) (3.546,2.211) (3.651,2.454) (3.756,2.699) (3.860,2.945) (3.965,3.190) (4.070,3.436) (4.174,3.680) (4.279,3.923) (4.384,4.163) (4.488,4.400) (4.593,4.634) (4.698,4.863) (4.802,5.088) (4.907,5.308) (5.011,5.521) (5.116,5.729) (5.221,5.930) (5.325,6.123) (5.430,6.309) (5.535,6.487) (5.639,6.656) (5.744,6.816) (5.849,6.967) (5.953,7.109) (6.058,7.241) (6.163,7.363) (6.267,7.475) (6.372,7.576) (6.476,7.666) (6.581,7.746) (6.686,7.815) (6.790,7.873) (6.895,7.920) (7.000,7.957) (7.104,7.982) (7.209,7.996) (7.314,8.000) (7.418,7.993) (7.523,7.975) (7.628,7.947) (7.732,7.909) (7.837,7.861) (7.941,7.802) (8.046,7.735) (8.151,7.658) (8.255,7.572) (8.360,7.478) (8.465,7.375) (8.569,7.264) (8.674,7.146) (8.779,7.021) (8.883,6.889) (8.988,6.750) (9.093,6.606) (9.197,6.456) (9.302,6.301) (9.406,6.141) (9.511,5.978) (9.616,5.810) (9.720,5.640) (9.825,5.466) (9.930,5.291) (10.034,5.113) (10.139,4.934) (10.244,4.754) (10.348,4.574) (10.453,4.394) (10.558,4.214) (10.662,4.034) (10.767,3.856) (10.871,3.680) (10.976,3.505) (11.081,3.333) (11.185,3.163) (11.290,2.996) (11.395,2.833) (11.499,2.673) (11.604,2.517) (11.709,2.365) (11.813,2.217) (11.918,2.074) (12.023,1.935) (12.127,1.802) (12.232,1.673) (12.336,1.549) (12.441,1.431) (12.546,1.318) (12.650,1.211) (12.755,1.108) (12.860,1.011) (12.964,0.920) (13.069,0.834) (13.174,0.753) (13.278,0.677) (13.383,0.607) (13.488,0.541) (13.592,0.481) (13.697,0.425) (13.802,0.374) (13.906,0.327) (14.011,0.284) (14.115,0.245) (14.220,0.211) (14.325,0.180) (14.429,0.152) (14.534,0.128) (14.639,0.106) (14.743,0.087) (14.848,0.071) (14.953,0.057) (15.057,0.046) (15.162,0.036) (15.267,0.028) (15.371,0.021) (15.476,0.016) (15.580,0.011) (15.685,0.008) (15.790,0.006) (15.894,0.004) (15.999,0.002) (16.104,0.001) (16.208,0.001) (16.313,0.000) (16.418,0.000) (16.522,0.000) (16.627,0.000) (16.732,0.000) (16.836,0.000) (16.941,0.000) (17.045,0.000)};
      % riga 0nubb a Q_bb
      \draw[zeroNu, line width=1.4pt] plot[domain=17.045-0.35:17.045+0.35, samples=90, smooth]
        (\x,{5.2*exp(-((\x-17.045)^2)/(2*0.045^2))});
      % asse
      \draw[white, opacity=0.7, line width=0.5pt] (2.5,0) -- (18.5,0);
      \draw[white, opacity=0.7, line width=0.5pt] (17.045,0) -- (17.045,-0.18);
      \node[anchor=north, white, opacity=0.85] at (17.045,-0.25) {\small$Q_{\beta\beta}$};
      \node[anchor=south west, twoNu] at (10.3,4.7) {\small$2\nu\beta\beta$};
      \node[anchor=south, zeroNu] at (17.045,5.35) {\small$0\nu\beta\beta$};
      \node[anchor=north west, inner sep=0pt, text width=16cm, align=left, white, opacity=0.75] at (2.5,-1.1)
        {\footnotesize\itshape Summed energy of the two electrons in ${}^{76}$Ge double beta decay, $Q_{\beta\beta} = 2039$~keV.};
    \end{scope}

    % titolo
    \node[anchor=north west, text width=16cm, inner sep=0pt, align=left]
      at (2.5,25.2)
      {\fontsize{46}{52}\selectfont\bfseries\spaceskip=0.3em plus 0.1em\relax\color{white} Neutrino Physics\par
       \vspace{14pt}
       {\fontsize{17}{22}\selectfont\spaceskip=0pt\relax\mdseries\itshape\color{twoNu}
        Oscillations, masses and double beta decay\par}};

    % autore
    \node[anchor=south west, inner sep=0pt] at (2.5,1.8)
      {\fontsize{15}{18}\selectfont\bfseries\color{white} Benedetta Rasera};
  \end{scope}
\end{tikzpicture}
\end{titlepage}